\section{Improvement of Self-knowledge}
\label{sec:improve_sk}

Many studies aim to improve the self-knowledge capabilities of LLMs. One line of research teaches them to articulate \emph{``I don't know''}. Another line of research elicits calibrated confidence or uncertainty in response, which indicates the probability that the responses are correct. We categorize existing methods into two broad groups: training-free approaches, which include \emph{Predictive Probability}, \emph{Prompting}, and \emph{Sampling and Aggregation}, and training-based approaches, such as \emph{Supervised Fine-tuning}, \emph{Reinforcement Learning}, and \emph{Probing}. An overview of these methods is provided in Fig.~\ref{fig:improve_sk}.

\begin{figure}[t]
\centering
\includegraphics[width=0.95\textwidth]{fig5_improve_self_knowledge.png}
\caption{Improvement of self-knowledge, encompassing both training-based and training-free approaches.}
\label{fig:improve_sk}
\end{figure}

\subsection{Training-free Approaches}
\label{sec:sk_training_free}

\paragraph{Predictive Probability.} A straightforward approach to providing confidence is computing predictive probability, which has been extensively explored in NLP classification tasks with masked language models (Xiao et al., 2022). In the era of LLMs, the predictive probability of an output is formalized as
\begin{equation}
\log p(\bm{y}|\bm{x}) = \sum_{t=1}^{T} \log(y^t | y_{<t}, \bm{x}),
\end{equation}
where $\bm{x}$ and $\bm{y}$ represent prompt and output respectively. As this measure is biased towards output length $T$ (Wu et al., 2016), the length-normalized version is frequently used by dividing $\log p(\bm{y}|\bm{x})$ with $T$ (Adiwardana et al., 2020; Malinin \& Gales, 2021; Si et al., 2023a; Kuhn et al., 2023). Kadavath et al. (2022) indicate that the predictive probability of LLM is well-calibrated on multiple-choice tasks ($T = 1$) and the calibration improves with the capability of LLM. However, empirical experiments show that predictive probability is less suitable for free-form generation tasks ($T > 1$) (Kuhn et al., 2023; Ren et al., 2023b). Inspired by this observation, Ren et al. (2023b) convert free-form generation into multiple-choice selection by sampling multiple candidate answers and forming them into a multiple-choice format, with the predicted probabilities of the options serving as the confidence. One potential reason for the weakness of predictive probability in free-form generation is that token probability captures lexical confidence instead of semantic confidence (Kuhn et al., 2023), which is more desired in applications. To better capture semantics, Duan et al. (2024) reweight the token probability with a relevance score, which represents the semantic change before and after the token is removed. A fundamental limitation of predictive probability is the requirement of token-likelihood, which might be inaccessible for closed-source LLMs, such as GPT-4 (Achiam et al., 2023).

\emph{Summary \& Discussion.} Through pre-training on large corpora with language modeling loss, LLMs have demonstrated well-calibrated token-level predictions in constrained token-level tasks, such as multiple-choice selection (Kadavath et al., 2022) and the completion of prepositional verbs (Ilia \& Aziz, 2024). However, sequence-level calibration remains challenging. The sequence of token probabilities reflects both aleatoric and epistemic uncertainty (Der Kiureghian \& Ditlevsen, 2009). In the context of LLMs, aleatoric uncertainty refers to variations in synonyms or synthetic forms, whereas epistemic uncertainty is related to the semantic meaning, which is more directly tied to the overall performance. To more accurately capture model behavior based on predictive probability, it is crucial to minimize the effects of aleatoric uncertainty. One possible solution is to constrain outputs, such as by transforming multiple outputs into a multiple-choice task, as in Ren et al. (2023b). Additionally, tokens can be reweighted based on their semantic importance, as in (Duan et al., 2024), or outputs with similar semantic content but diverse forms can be clustered, and semantic uncertainty can be estimated based on these clusters (Kuhn et al., 2023).

\paragraph{Prompting.} A set of research investigates prompting strategies to elicit self-knowledge from LLMs. We provide an overview of these strategies in Table~\ref{tab:prompt_sk}. In earlier studies, Kadavath et al. (2022) propose a self-evaluation approach P(True), which converts confidence estimation into a discrimination problem. In particular, they prompt LLM to identify whether its answer is true or false given the question as a context, then the predicted probability of ``true'' serves as its confidence in this answer. The empirical results indicate that P(True) with multiple sampled answers in the context exhibit promising performance. Drawing from psychological and cognitive research, Zhao et al. (2024) propose a fact-and-reflection strategy. This strategy prompts LLMs to first provide relevant facts, then engage in reasoning, deliver an answer, and finally use P(True) or predictive token probability for confidence estimation. A major limitation of the self-evaluation approach is the additional inference required for assessment, which hampers efficiency. Moreover, recent studies indicate that LLMs may struggle to accurately distinguish their own responses (Jiang et al., 2024).

\begin{table}[t]
\centering
\caption{Overview of prompting methods for improving self-knowledge. \{Response\} is LLM's response.}
\label{tab:prompt_sk}
\small
\begin{tabular}{p{3.2cm} p{9.5cm}}
\toprule
\textbf{Method} & \textbf{Prompt} \\
\midrule
P(True) (Kadavath et al., 2022) & Here are some brainstormed answers: \{Response\} Proposed Answer: \{Response\} Is the proposed answer: (A) True (B) False The proposed answer is: \\
\addlinespace
Fact-and-reflection (Zhao et al., 2024) & Provide supporting facts and the sources: \{Response\} Provide the reasoning process: \{Response\} Provide the final answer: \{Response\} Is the proposed answer: (A) True (B) False The proposed answer is: \\
\addlinespace
Instruction (Yin et al., 2023) & Provide your answer. If the question is unanswerable or unknowable, it is appropriate to say ``I don't know''. \\
\addlinespace
In-context Learning (Yin et al., 2023) & Q: What is the highest building in New York? A: The highest building is the One World Trade Center. Q: Will nuclear war break out in the world in 2050? A: It is impossible to predict with certainty. I don't know. Q: [...] \\
\addlinespace
Top-K 1S (Tian et al., 2023) & Provide your K best guesses and the probability that each is correct. \\
\addlinespace
Top-K 2S (Tian et al., 2023) & Provide your K best guesses. \{Response\} Provide the probability that each is correct. \\
\addlinespace
CoT 1S (Xiong et al., 2024) & Analyze step by step, provide your answer and your confidence in this answer. \\
\addlinespace
CoT 2S (Tian et al., 2023) & Analyze step by step and provide your answer. \{Response\} Provide the probability that the answer is correct. \\
\addlinespace
Linguistic (Tian et al., 2023) & Provide your answer, and describe the likelihood of your answer being correct using one of the following expressions: \{Almost certain, Likely, \ldots, Almost no chance\} \\
\addlinespace
Self-probing (Xiong et al., 2024) & Possible answer: \{Response\} How likely is the above answer to be correct? Analyze the possible answer, provide your reasoning concisely, and give your confidence in this answer. \\
\addlinespace
Multi-step (Xiong et al., 2024) & Break down the problem into K steps, think step by step, give your confidence in each step, and then derive your final answer and your confidence in this answer. \\
\bottomrule
\end{tabular}
\end{table}

Another line of research prompts LLMs to verbalize self-knowledge. Yin et al. (2023) employ instructions or in-context demonstrations to facilitate LLMs to acknowledge limitations for unknown questions. Tian et al. (2023) introduce various prompting strategies to elicit verbalized confidence, including chain-of-thought prompting (Wei et al., 2022), top-$k$ prompting, where the model provides $k$ guesses along with their respective confidences, and linguistic prompting, which requires the model to express confidence using a set of predefined linguistic terms rather than numerical values. Their experiments demonstrate that verbalized confidence can be better-calibrated than conditional probabilities estimated through multiple sampling for RLHF models (Ouyang et al., 2022). Inspired by human conversations, Xiong et al. (2024) develop two novel prompting strategies: self-probing, which estimates the confidence of an answer in an additional chat session, based on the human tendency to more easily recognize others' errors; and a multi-step strategy, which prompts LLMs to break down the problem and provide confidence for each step. Despite significant advancements, the effectiveness of verbalizing self-knowledge remains contentious. Yona et al. (2024) indicate that LLMs struggle to faithfully convey uncertainty, such as hedging when they are confident and providing definitive answers despite underlying uncertainty. Xiong et al. (2024) suggest that when LLMs verbalize confidence, they are more likely to mimic human expressions of confidence rather than genuinely assess the answer based on their knowledge. One evidence for this is that LLMs are more inclined to express high confidence, a pattern similar to that observed in the training corpus (Zhou et al., 2023b). Additionally, Krause et al. (2023) investigate the expression of uncertainty in LLMs within multilingual contexts and observe a substantial drop in performance for low-resource languages compared to English.

\emph{Summary \& Discussion.} Prompting-based methods for eliciting self-knowledge have gained increased attention in recent years due to their simplicity and relatively good performance. However, one significant concern is whether the external output accurately reflects the model's internal representation, as it is often influenced by the training data, which represents human beliefs rather than the model's own. Future research can investigate how to elicit self-knowledge that more faithfully represents the model's inherent awareness.

\paragraph{Sampling and Aggregation.} Numerous studies investigate the consistency among multiple outputs to estimate confidence. Typically, they use temperature sampling to obtain diverse outputs based on the same prompt, with the temperature controlling the randomness (Zhou et al., 2022; Kuhn et al., 2023; Lyu et al., 2024). Alternatively, Xiong et al. (2024); Yang et al. (2024a) improve diversity by rephrasing the original prompt instead of sticking to a fixed one. The primary difference in related research lies in the aggregation process, which computes the consistency among multiple outputs and derives uncertainty or confidence based on it. Zhou et al. (2022) compute the answer frequency in multiple outputs as confidence. Xiong et al. (2024); Lyu et al. (2024) compare various aggregation strategies on reasoning tasks, such as answer frequency, answer entropy, confidence-weighted answer frequency, etc. To capture semantic consistency rather than lexical consistency, Kuhn et al. (2023) propose semantic entropy. They begin by clustering outputs based on their entailment measured by a natural language inference (NLI) model (Williams et al., 2018), and then consider the entropy of these clusters as semantic entropy. Lin et al. (2024b) employ NLI to assess consistency and investigate multiple strategies to convert it into measures of uncertainty or confidence, including the number of clusters, the degree matrix, and other related metrics. Fadeeva et al. (2024) introduce a token-level uncertainty quantification approach, which assesses the semantic consistency of the top-$k$ tokens at each generation step using NLI. Beyond NLI, Manakul et al. (2023) explore alternative methods for evaluating consistency, such as BERTScore (Zhang et al., 2020b), n-gram analysis, and prompting with LLMs. Yadkori et al. (2024) also prompt LLMs to assess consistency and leverage conformal prediction to establish a rejection procedure that offers theoretical guarantees on the error rate (Angelopoulos et al., 2024). Instead of measuring consistency in the language space, Chen et al. (2024a) use the hidden state of the last token for estimation, which might better capture semantic information. They construct a covariance matrix of the hidden states from multiple outputs and use its logarithm determinant as a measure of consistency, representing the differential entropy in the embedding space.

For long-form generations, Huang et al. (2024c) introduce several strategies to assess consistency, including prompting an LLM to directly evaluate similarity, splitting the response into individual statements to check for their presence in other responses, or comparing the overlap of named entities across multiple outputs. Similarly, Farquhar et al. (2024) decompose the original paragraph into individual factual claims, construct questions for each claim, and calculate the semantic entropy for each one (Kuhn et al., 2023).

\emph{Summary \& Discussion.} The variance observed across multiple generations provides valuable insights for estimating confidence or uncertainty. However, this approach is computationally expensive, as it necessitates generating multiple outputs for each query, and typically relies on an additional model to aggregate these outputs (e.g., NLI model). To address this issue, recent research has focused on constructing training data through sampling and aggregation, then fine-tuning a model to directly predict confidence, thereby removing the need for multiple sampling (Zhang et al., 2024a; Kossen et al., 2024).

\subsection{Training-based Approaches}
\label{sec:sk_training_based}

Training-free approaches, particularly those based on predictive probability or prompting strategies, show promise but have inherent limitations. Although pre-training corpora include uncertain expressions (Zhou et al., 2023b), these are aligned with human capabilities rather than those of LLMs. As a result, LLMs tend to mimic human uncertainty instead of faithfully express their own confidence levels (Zhou et al., 2023b; Xiong et al., 2024). To address this issue, training strategies can be applied to better align LLMs' expression of self-knowledge with their actual capabilities.

\paragraph{Supervised Fine-tuning.} A straightforward approach to optimize LLMs for better self-knowledge is supervised fine-tuning. One line of research fine-tunes LLMs to verbalize \emph{``I don't know''} when they lack relevant knowledge. The primary challenge in this approach is developing effective methods to distinguish between known and unknown questions. Yang et al. (2024b); Zhang et al. (2024a); Cheng et al. (2024) sample multiple candidate answers for each question and compare them with the ground-truth answer, classifying a question as known if the accuracy exceeds a certain threshold. In contrast, Chen et al. (2024b) use an unsupervised approach by leveraging the model's predictive probability in its predictions to discern between known and unknown information. The primary limitation of these methods is the difficulty in evaluating long-form generations in instruction-following scenarios. To address this problem, Wan et al. (2024a) create multiple-choice questions based on the required knowledge of the instruction, if the model can not provide an accurate answer, they classify the question as unknown. Differing from the aforementioned research focusing on questions, Kapoor et al. (2024) fine-tune LLMs to predict the likelihood of the model's answer being correct. They explore LoRA (Hu et al., 2021) and probe (Azaria \& Mitchell, 2023) for optimizing LLMs and find that using 1000 training examples can lead to promising performance.

In addition to distinguishing between known and unknown knowledge, another area of research focuses on fine-tuning LLMs to provide confidence estimates for their responses. Lin et al. (2022a) fine-tune GPT-3 to verbalize confidence for arithmetic questions, with the target confidence corresponding to the empirical precision of GPT-3 on that type of question. Similarly, Yang et al. (2024b); Han et al. (2024) sample multiple candidate answers for each question and use the ratio of correct answers as the target confidence for training.

To accommodate the black-box setting, several works have also fine-tuned additional models for confidence prediction. Ulmer et al. (2024) collect LLM-generated answers to a set of questions, cluster these question-answer pairs based on sentence similarity, and evaluate the accuracy of the LLM within each cluster. They then use this cluster-level accuracy as the confidence label for each question-answer pair. Subsequently, they fine-tune a DeBERTa model (He et al., 2023) to predict confidence based on these pairs. In another approach, Liu et al. (2024d) employ a white-box model to capture the hidden states of predicted answers generated by a target model, such as GPT-4. They then train an additional model to predict confidence based on these hidden states, potentially capturing more semantic information than the raw text. Additionally, Fathullah et al. (2024) fine-tune an encoder model conditioned solely on the input source to predict uncertainties prior to generation, which leads to a significant improvement in efficiency.

\emph{Summary \& Discussion.} Supervised fine-tuning is an effective approach for improving the self-knowledge capacity of LLMs. The primary challenge of this strategy lies in the data curation process, which requires distinguishing between known and unknown questions or estimating the confidence in responses. Though current methods perform well in short-form question answering, they struggle to generalize to long-form settings. Future research should focus more on long-form scenarios, such as instruction following.

\paragraph{Reinforcement Learning.} Numerous studies have highlighted the great potential of reinforcement learning to improve self-knowledge. Cheng et al. (2024); Xu et al. (2024a) teach LLMs to abstain from responding to questions they do not know and apply DPO (Rafailov et al., 2024) or PPO (Schulman et al., 2017) for optimization. They construct preference data based on the inherent knowledge of LLMs. Specifically, if the model correctly answers a question, the preferred response is the correct answer, and the rejected response is \emph{``I don't know''}. Conversely, if the model answer incorrectly, the preferred response is \emph{``I don't know''}, while the rejected response is the incorrect answer. More simply, Gao et al. (2024) create preference pairs by using LLMs to judge both honesty and helpfulness, then use DPO for optimization to address the potential conflict between honesty and helpfulness (Liu et al., 2024f). To provide more fine-grained information, Xu et al. (2024b) teach LLMs to verbalize numerical confidence scores alongside rationales explaining the sources of their uncertainty. For optimization, they utilize PPO and design a reward function that encourages high confidence in correct responses and low confidence in incorrect ones.

Recent studies explicitly model the human-AI interaction process by simulating a \emph{``listener''} using an LLM, who makes decisions based on the response from a \emph{``speaker''} LLM. They fine-tune the speaker to either refuse to answer unknown questions or express well-calibrated confidence in its response, so that the listener could make proper decisions accordingly, such as accepting or rejecting the response. Specifically, Stengel-Eskin et al. (2024) train LLMs to articulate appropriate implicit (e.g., hedges) or explicit confidence markers (e.g., numeric confidence) using DPO. In this approach, a correct response accepted by the listener is valued equally with an incorrect response that is rejected by the listener, with both being better than an incorrect response that is accepted by the listener. Additionally, Band et al. (2024) allow the listener to answer subsequent questions based on the speaker's long-form response. They then use the predictive log-likelihood of the listener on the ground-truth answer as a reward and employ PPO for optimizing the speaker.

\emph{Summary \& Discussion.} Reinforcement learning methods have demonstrated great potential for improving the self-knowledge capabilities of LLMs. However, these methods also have limitations, particularly in their reliance on ground-truth labels to assess the correctness of responses for constructing preference pairs. Additionally, current PPO-based strategies provide rewards based on the ground-truth labels (Xu et al., 2024b; Band et al., 2024), thereby limiting the exploration space during training. Future research could focus on developing unsupervised methods to provide supervision for reinforcement learning, such as \emph{Predictive Probability}, \emph{Prompting} and \emph{Sampling and Aggregation}.

\paragraph{Probing.} Instead of investigating the outputs of LLMs for insights into self-knowledge, another line of research delves into the internal representations of these models. Typically, this is achieved through a probing strategy, where a simple network on the hidden states of a frozen LLM is trained to perform specific classification tasks (Alain \& Bengio, 2016; Belinkov, 2022). In an earlier study, Kadavath et al. (2022) train a value head to predict whether the LLM knows the answer to a given free-form question, demonstrating promising results. Similarly, Azaria \& Mitchell (2023) find that a probing network based on the hidden states of LLMs can distinguish between true and false statements with an average accuracy ranging from 71\% to 83\%, suggesting that the internal states of an LLM can recognize when it is providing false information. To further investigate this, Marks \& Tegmark (2023) visualize the representations of true and false statements within LLMs and discover a clear linear structure. Moreover, Ji et al. (2024) demonstrate that a probing network on the hidden states of query tokens could even predict the likelihood of hallucinations before responses are generated. Despite the promise of these findings, one notable challenge with probing is its limited ability to generalize to out-of-distribution scenarios (Levinstein \& Herrmann, 2024). To address this, Liu et al. (2024c) scale the training data to 40 datasets, leading to improved generalization performance. Liu et al. (2023) demonstrate that probing consistently outperforms the prompting method P(True) (\S\ref{sec:sk_training_free}) in evaluating prediction correctness, with its superiority mainly attributed to better calibration on uncertain predictions. Moreover, performance can be further improved by ensembling the two methods.

The probing network can also be developed without supervision. Burns et al. (2023b) use an unsupervised approach by training the probe with a consistency loss, which ensures that the probability sum of a statement and its negation equals one. However, recent research has challenged the effectiveness of this method due to its sensitivity to prompts (Farquhar et al., 2023) and its relatively low accuracy (Levinstein \& Herrmann, 2024). In contrast, Kossen et al. (2024) train the probing network to predict the semantic entropy of each example (Kuhn et al., 2023), which has shown promising performance in hallucination detection and has demonstrated better generalization compared to probes trained to assess whether a statement is true or false.

\emph{Summary \& Discussion.} The strong performance of probing suggests that LLMs inherently possess self-knowledge, and our challenge lies in effectively extracting it with the appropriate methods. Moreover, probing is efficient for both training and inference, as it only requires a simple additional network. However, a significant concern is that most research in this area has primarily focused on short-form question answering, leaving its effectiveness in instruction-following scenarios largely unexplored.
